% GENERATED FILE. Edit canonical lessons, then run npm run generate:books.
% canonical-source-hash: fnv1a64:0c3b6bdb9b888ab1
% canonical-lessons: GU-C518-pasti, GU-C518-jalu, GU-C518-jhadu, GU-C518-phug, GU-C518-cotadvu

\chapter{Clearing Out the Loft}
\label{ch:gu-clearing-out-the-loft}
\hlchaptermodality{518}

\begin{chapteropening}
\textbf{By the end of this chapter:} \emph{I can talk about clearing out old paper, cobwebs and mould.}

Complete the last lesson of chapter 518: I can talk about clearing out old paper, cobwebs and mould.
\end{chapteropening}

\section[\texorpdfstring{\gu{પસ્તી} (pastī)}{પસ્તી (pastī)}]{\gu{પસ્તી} (pastī) — waste paper (sold by weight)}

\label{lesson:GU-C518-pasti}



\begin{quote}
Before the new one: say the Gujarati for a light bulb, then the Gujarati for to push in (firmly).
\end{quote}

\subsection*{You'll want to know: \gu{પસ્તી}}
\textbf{\gu{પસ્તી}} — \emph{pastī} — \textquotedblleft{}waste paper (sold by weight)\textquotedblright{}.

\begin{quote}
\gu{અમે} \gu{મહિને} \gu{એક} \gu{વાર} \gu{પસ્તી} \gu{વેચીએ} \gu{છીએ}.
\end{quote}

\begin{quote}
\emph{ame mahine ek vār pastī vecīe chīe.}
\end{quote}

\begin{quote}
Once a month we sell our old newspapers.
\end{quote}

A feminine noun.

\begin{cousinweb}
\textbf{\gu{પસ્તી}} (\emph{pastī}) is old newspaper and paper, sold to a dealer by weight. No dictionary we use records an older form, so its history is not settled.
\end{cousinweb}

\begin{grammarlens}[title={What you've built}]
One more word for describing what you lived and felt.
\end{grammarlens}

\subsection*{Your turn}
\begin{itemize}
\raggedright
  \item \emph{Say it:} \emph{pastī}
  \item \emph{Say it:} \emph{pastī}, once more
  \item \emph{Say it:} \emph{ame mahine ek vār pastī vecīe chīe.} — Once a month we sell our old newspapers.
  \item \emph{Recall:} say the Gujarati for a light bulb, then the Gujarati for to push in (firmly), then say \emph{pastī} again
\end{itemize}

\begin{tcolorbox}[breakable,colback=teal!4,colframe=teal!35!black,title={Before you move on}]
What does \gu{પસ્તી} mean? (\textquotedblleft{}Waste paper (sold by weight)\textquotedblright{}.) Say it once more.
\end{tcolorbox}
\section[\texorpdfstring{\gu{જાળું} (jāḷũ)}{જાળું (jāḷũ)}]{\gu{જાળું} (jāḷũ) — a cobweb}

\label{lesson:GU-C518-jalu}



\begin{quote}
Before the new one: say the Gujarati for to push in (firmly), then the Gujarati for waste paper.
\end{quote}

\subsection*{You'll want to know: \gu{જાળું}}
\textbf{\gu{જાળું}} — \emph{jāḷũ} — \textquotedblleft{}a cobweb\textquotedblright{}.

\begin{quote}
\gu{છતના} \gu{ખૂણામાં} \gu{કરોળિયાનું} \gu{જાળું} \gu{છે}.
\end{quote}

\begin{quote}
\emph{chatnā khūṇāmā̃ karoḷiyānũ jāḷũ che.}
\end{quote}

\begin{quote}
There is a spider's web in the corner of the ceiling.
\end{quote}

A neuter noun.

\begin{cousinweb}
\textbf{\gu{જાળું}} (\emph{jāḷũ}) is inherited from Sanskrit \emph{jāla}, \textquotedblleft{}a net\textquotedblright{}, and CDIAL lists the Gujarati word, \textquotedblleft{}a net, a cobweb\textquotedblright{} (CDIAL 5213).
\end{cousinweb}

\begin{grammarlens}[title={What you've built}]
One more word for describing what you lived and felt.
\end{grammarlens}

\subsection*{Your turn}
\begin{itemize}
\raggedright
  \item \emph{Say it:} \emph{jāḷũ}
  \item \emph{Say it:} \emph{jāḷũ}, once more
  \item \emph{Say it:} \emph{chatnā khūṇāmā̃ karoḷiyānũ jāḷũ che.} — There is a spider's web in the corner of the ceiling.
  \item \emph{Recall:} say the Gujarati for to push in (firmly), then the Gujarati for waste paper, then say \emph{jāḷũ} again
\end{itemize}

\begin{tcolorbox}[breakable,colback=teal!4,colframe=teal!35!black,title={Before you move on}]
What does \gu{જાળું} mean? (\textquotedblleft{}A cobweb\textquotedblright{}.) Say it once more.
\end{tcolorbox}
\section[\texorpdfstring{\gu{ઝાડુ} (jhāḍu)}{ઝાડુ (jhāḍu)}]{\gu{ઝાડુ} (jhāḍu) — a long broom}

\label{lesson:GU-C518-jhadu}



\begin{quote}
Before the new one: say the Gujarati for waste paper, then the Gujarati for a cobweb.
\end{quote}

\subsection*{You'll want to know: \gu{ઝાડુ}}
\textbf{\gu{ઝાડુ}} — \emph{jhāḍu} — \textquotedblleft{}a long broom\textquotedblright{}.

\begin{quote}
\gu{મીરાએ} \gu{ઝાડુથી} \gu{આખો} \gu{ઓરડો} \gu{સાફ} \gu{કર્યો}.
\end{quote}

\begin{quote}
\emph{Mīrāe jhāḍuthī ākho orḍo sāph karyo.}
\end{quote}

\begin{quote}
Mira swept the whole room clean with the broom.
\end{quote}

A neuter noun.

\begin{cousinweb}
\textbf{\gu{ઝાડુ}} (\emph{jhāḍu}) is perhaps related to \emph{jhāṛvũ}, \textquotedblleft{}to shake, to sift\textquotedblright{}, which CDIAL files under a reconstructed \emph{*jhaṭati} with Panjabi \emph{jhāṛṇā}, \textquotedblleft{}to shake off, to sweep\textquotedblright{} (CDIAL 5328); CDIAL does not list the broom itself, so its history is not settled.
\end{cousinweb}

\begin{grammarlens}[title={What you've built}]
One more word for describing what you lived and felt.
\end{grammarlens}

\subsection*{Your turn}
\begin{itemize}
\raggedright
  \item \emph{Say it:} \emph{jhāḍu}
  \item \emph{Say it:} \emph{jhāḍu}, once more
  \item \emph{Say it:} \emph{Mīrāe jhāḍuthī ākho orḍo sāph karyo.} — Mira swept the whole room clean with the broom.
  \item \emph{Recall:} say the Gujarati for waste paper, then the Gujarati for a cobweb, then say \emph{jhāḍu} again
\end{itemize}

\begin{tcolorbox}[breakable,colback=teal!4,colframe=teal!35!black,title={Before you move on}]
What does \gu{ઝાડુ} mean? (\textquotedblleft{}A long broom\textquotedblright{}.) Say it once more.
\end{tcolorbox}
\section[\texorpdfstring{\gu{ફૂગ} (phūg)}{ફૂગ (phūg)}]{\gu{ફૂગ} (phūg) — mould (a fungus)}

\label{lesson:GU-C518-phug}



\begin{quote}
Before the new one: say the Gujarati for a cobweb, then the Gujarati for a long broom.
\end{quote}

\subsection*{You'll want to know: \gu{ફૂગ}}
\textbf{\gu{ફૂગ}} — \emph{phūg} — \textquotedblleft{}mould (a fungus)\textquotedblright{}.

\begin{quote}
\gu{ભીની} \gu{દીવાલ} \gu{પર} \gu{ફૂગ} \gu{આવી} \gu{ગઈ} \gu{છે}.
\end{quote}

\begin{quote}
\emph{bhīnī dīvāl par phūg āvī gaī che.}
\end{quote}

\begin{quote}
Mould has appeared on the damp wall.
\end{quote}

A feminine noun.

\begin{cousinweb}
\textbf{\gu{ફૂગ}} (\emph{phūg}) grows where it is damp; it \textquotedblleft{}comes\textquotedblright{} on a thing: \textbf{\gu{ફૂગ} \gu{આવી} \gu{ગઈ}}. No dictionary we use records an older form, so its history is not settled.
\end{cousinweb}

\begin{grammarlens}[title={What you've built}]
One more word for describing what you lived and felt.
\end{grammarlens}

\subsection*{Your turn}
\begin{itemize}
\raggedright
  \item \emph{Say it:} \emph{phūg}
  \item \emph{Say it:} \emph{phūg}, once more
  \item \emph{Say it:} \emph{bhīnī dīvāl par phūg āvī gaī che.} — Mould has appeared on the damp wall.
  \item \emph{Recall:} say the Gujarati for a cobweb, then the Gujarati for a long broom, then say \emph{phūg} again
\end{itemize}

\begin{tcolorbox}[breakable,colback=teal!4,colframe=teal!35!black,title={Before you move on}]
What does \gu{ફૂગ} mean? (\textquotedblleft{}Mould (a fungus)\textquotedblright{}.) Say it once more.
\end{tcolorbox}
\section[\texorpdfstring{\gu{ચોંટાડવું} (cõṭāḍvũ)}{ચોંટાડવું (cõṭāḍvũ)}]{\gu{ચોંટાડવું} (cõṭāḍvũ) — to stick (something) on}

\label{lesson:GU-C518-cotadvu}



\begin{quote}
Before the new one: say the Gujarati for a long broom, then the Gujarati for mould.
\end{quote}

\subsection*{You'll want to know: \gu{ચોંટાડવું}}
\textbf{\gu{ચોંટાડવું}} — \emph{cõṭāḍvũ} — \textquotedblleft{}to stick (something) on\textquotedblright{}.

\begin{quote}
\gu{મીરાએ} \gu{પેટી} \gu{પર} \gu{કાગળ} \gu{ચોંટાડ્યો}.
\end{quote}

\begin{quote}
\emph{Mīrāe peṭī par kāgaḷ cõṭāḍyo.}
\end{quote}

\begin{quote}
Mira stuck a piece of paper on the box.
\end{quote}

Past, with \gu{તેણે} (he or she …): \gu{તેણે} \gu{કાગળ} \gu{ચોંટાડ્યો}, \emph{teṇe kāgaḷ cõṭāḍyo}.

\begin{cousinweb}
\textbf{\gu{ચોંટાડવું}} (\emph{cõṭāḍvũ}) is built from \textbf{\gu{ચોંટવું}}, \textquotedblleft{}to stick\textquotedblright{}, with the causative \emph{-āḍ-}: to make something stick. The history of \textbf{\gu{ચોંટવું}} itself is not settled.
\end{cousinweb}

\begin{grammarlens}[title={What you've built}]
One more word for describing what you lived and felt.
\end{grammarlens}

\subsection*{Your turn}
\begin{itemize}
\raggedright
  \item \emph{Say it:} \emph{cõṭāḍvũ}
  \item \emph{Say it:} \emph{cõṭāḍvũ}, once more
  \item \emph{Say it:} \emph{Mīrāe peṭī par kāgaḷ cõṭāḍyo.} — Mira stuck a piece of paper on the box.
  \item \emph{Recall:} say the Gujarati for a long broom, then the Gujarati for mould, then say \emph{cõṭāḍvũ} again
\end{itemize}

\begin{tcolorbox}[breakable,colback=teal!4,colframe=teal!35!black,title={Before you move on}]
What does \gu{ચોંટાડવું} mean? (\textquotedblleft{}To stick (something) on\textquotedblright{}.) Say it once more.
\end{tcolorbox}
